% Authored mathematical statements for the fixed E2 computation construction and proof groups.

\begin{definition}[The detector target in degree (19,114)]
\label{frontier:E2Computation:0c823d7609b2}
\level{fine}
Fix the selected element $B^{\mathrm{src}}_{19,114}\in D^{19,114}$ of the detector source component.
\lean{KIP126.Computation.Standard.detectorBetaGFour}
\end{definition}

\begin{definition}[The bounded sphere comparison]
\label{frontier:E2Computation:1359a65ba983}
\level{fine}
For $s,t\in\mathbb N$ with $t\leq261$, let
\[c_{s,t}:P^{s,t}\cong_{\mathbb Z}E_2^{s,t}(S^0)\]
be the fixed linear equivalence.
\lean{KIP126.Computation.Standard.comparison}
\end{definition}

\begin{definition}[The detector source components]
\label{frontier:E2Computation:1dfdcb9f6b7a}
\level{fine}
For every $s,t\in\mathbb N$, let $D^{s,t}$ be the fixed $\mathbb F_2$-module of detector source classes in bidegree $(s,t)$.
\lean{KIP126.Computation.Standard.detectorComponent}
\end{definition}

\begin{definition}[The bounded detector comparison]
\label{frontier:E2Computation:3be7c427b9bf}
\level{fine}
For $s,t\in\mathbb N$ with $t\leq261$, let
\[c^D_{s,t}:D^{s,t}\cong_{\mathbb Z}E_2^{s,t}(D)\]
be the fixed linear equivalence, where $D$ is the detector spectrum.
\lean{KIP126.Computation.Standard.detectorComparison}
\end{definition}

\begin{definition}[Values of the source generators]
\label{frontier:E2Computation:5978caa8605f}
\level{fine}
For every index $i<N$ in the fixed generator catalogue, with $\deg(i)=(s_i,t_i)$, fix the source generator value $u_i\in P^{s_i,t_i}$.
\lean{KIP126.Computation.Standard.sourceGenerator}
\end{definition}

\begin{definition}[The sphere source components]
\label{frontier:E2Computation:5e3e82d48166}
\level{fine}
For every $s,t\in\mathbb N$, let $P^{s,t}$ be the fixed $\mathbb F_2$-module of sphere source classes in bidegree $(s,t)$.
\lean{KIP126.Computation.Standard.sourceComponent}
\end{definition}

\begin{definition}[Indices of the named atoms]
\label{frontier:E2Computation:5ff813cd73ee}
\level{fine}
For every named atom $a$, fix its index $\iota(a)<N$ in the sphere generator catalogue.
\lean{KIP126.Computation.Standard.atomIndex}
\end{definition}

\begin{definition}[The top-cell class of the nu-cofiber]
\label{frontier:E2Computation:6ac06ebae49b}
\level{fine}
Let $C\nu$ be the cofiber of the fixed map $\nu$. Fix the top-cell label $a\in E_2^{8,134}(C\nu)$.
\lean{KIP126.Computation.Standard.topCell}
\end{definition}

\begin{definition}[The sphere generator catalogue]
\label{frontier:E2Computation:7008d6a19504}
\level{fine}
Fix the finite sphere generator catalogue: a natural number $N$ and a record $(\operatorname{name}(i),s_i,t_i)$ for every $i<N$. Write $\deg(i)=(s_i,t_i)$.
\lean{KIP126.Computation.Standard.generators}
\end{definition}

\begin{definition}[Source multiplication, unit and named atoms]
\label{frontier:E2Computation:843c331404a2}
\level{fine}
The sphere source components are equipped with:
\begin{enumerate}
\item an element $a^{\mathrm{src}}\in P^{\deg(a)}$ for every named atom $a$;
\item an element $1_P\in P^{0,0}$;
\item for $s,t,s',t'\in\mathbb N$, a bilinear map
\[\mu_{s,t;s',t'}:P^{s,t}\times P^{s',t'}\longrightarrow P^{s+s',t+t'}.\]
\end{enumerate}
\lean{KIP126.Computation.Standard.sourceAtom,KIP126.Computation.Standard.sourceUnit,KIP126.Computation.Standard.sourceProduct}
\end{definition}

\begin{definition}[Coordinates in each bounded source degree]
\label{frontier:E2Computation:8f09e4b56647}
\level{fine}
For $s,t\in\mathbb N$ and $t\leq261$, let
\[q_{s,t}:P^{s,t}\cong_{\mathbb Z}\bigoplus_{i<n_{s,t}}\mathbb F_2\]
be the fixed coordinate equivalence, where $n_{s,t}$ is the archive basis count.
\lean{KIP126.Computation.Standard.sourceCoordinates}
\end{definition}

\begin{definition}[The archive basis counts]
\label{frontier:E2Computation:c04b5bfa5c79}
\level{fine}
For every $s,t\in\mathbb N$, fix the natural number $n_{s,t}$ of basis monomials in the sphere archive.
\lean{KIP126.Computation.Standard.basisCount}
\end{definition}

\begin{definition}[The archive basis monomials]
\label{frontier:E2Computation:ca5f76beefcc}
\level{fine}
For $s,t\in\mathbb N$ and $i<n_{s,t}$, let $m_{s,t,i}$ be the selected source monomial of bidegree $(s,t)$. A monomial is formed from the unit, catalogue generators and ordered products whose bidegrees add.
\lean{KIP126.Computation.Standard.basisMonomial}
\end{definition}

\begin{definition}[The detector class in degree (16,112)]
\label{frontier:E2Computation:cd42dbe2a65b}
\level{fine}
Fix the selected source element $v^{\mathrm{src}}_2{}^{16}\in D^{16,112}$.
\lean{KIP126.Computation.Standard.detectorV2Sixteen}
\end{definition}

\begin{definition}[The nu-cofiber coordinate vectors]
\label{frontier:E2Computation:762f8710f8d1}
\level{fine}
For every $s,t,i\in\mathbb N$, fix a coordinate vector
\[b_{C\nu;s,t,i}\in E_2^{s,t}(C\nu).\]
\lean{KIP126.Computation.Table.cofiberBasis}
\end{definition}

\begin{theorem}[Dimensions of the bounded sphere archive]
\label{frontier:E2Computation:75f3d8ca6b5b}
\level{fine}
For every $s,t\in\mathbb N$ with $t\leq261$, the archive basis count equals the tabulated sphere dimension:
\[n_{s,t}=\operatorname{dim}_{\mathrm{tab}}(s,t).\]
\lean{KIP126.Computation.Table.archive_dimension}
\end{theorem}

\begin{theorem}[Primitive product and obstruction relations]
\label{frontier:E2Computation:f5588f00ea38}
\level{fine}
The following statements hold; $G_{j;s,t}$ denotes the image of the addressed archive generator in the sphere $E_2$-page.
\begin{enumerate}
\item $h_2G_{407;13,135}=0$.
\item $h_0^6x_{125,8}=0$.
\item $h_0G_{24;5,36}=0$.
\item $h_0^5x_{125,9,2}=0$.
\item $h_5^2b=0$ in bidegree $(10,134)$.
\item $h_5G_{84;7,70}=0$.
\item $h_0^3G_{89;10,72}=0$.
\item $h_2Md_0=h_0G_{89;10,72}$.
\item $(h_0G_{3;1,8})G_{335;12,130}=h_0^2G_{425;12,137}$.
\item $h_5G_{85;7,70}=0$.
\end{enumerate}
\lean{KIP126.Computation.Table.h2_obstruction_rel0,KIP126.Computation.Table.h0_obstruction_rel3,KIP126.Computation.Table.h0_obstruction_rel0,KIP126.Computation.Table.h0_obstruction_rel2,KIP126.Computation.Standard.h5Square_b_zero,KIP126.Computation.Table.theta_h5_generator84_zero,KIP126.Computation.Table.h2_obstruction_rel2,KIP126.Computation.Table.h2_obstruction_rel1,KIP126.Computation.Table.h0_obstruction_rel1,KIP126.Computation.Table.theta_h5_generator85_zero}
\end{theorem}

\begin{theorem}[Calibration of the two detector labels]
\label{frontier:E2Computation:07bf2180c6a5}
\level{fine}
Let $\phi:\mathrm{tmf}\cong D$ be the fixed detector isomorphism, and let $v_2^{16}$ and $B$ be the selected literature classes in bidegrees $(16,112)$ and $(19,114)$ of the $\mathrm{tmf}$ Adams $E_2$-page. Then
\[c^D_{16,112}(v^{\mathrm{src}}_2{}^{16})=\phi_*(v_2^{16}),\qquad c^D_{19,114}(B^{\mathrm{src}}_{19,114})=\phi_*(B).\]
\lean{KIP126.Computation.Standard.detector_v2_agrees,KIP126.Computation.Standard.detector_beta_agrees}
\end{theorem}

\begin{theorem}[The comparison preserves the unit and bounded products]
\label{frontier:E2Computation:07da9aeedae4}
\level{fine}
The comparison satisfies $c_{0,0}(1_P)=1$. For $s,t,s',t'\in\mathbb N$, $t+t'\leq261$, $x\in P^{s,t}$ and $y\in P^{s',t'}$,
\[c_{s+s',t+t'}(\mu_{s,t;s',t'}(x,y))=c_{s,t}(x)c_{s',t'}(y).\]
\lean{KIP126.Computation.Standard.comparison_unit,KIP126.Computation.Standard.comparison_product}
\end{theorem}

\begin{theorem}[Exact identification of each atom record]
\label{frontier:E2Computation:2ab64b6c8429}
\level{fine}
For every named atom $a$, its selected index and catalogue record recover its prescribed archive record:
\[(\iota(a),\operatorname{name}(\iota(a)),s_{\iota(a)},t_{\iota(a)})=\operatorname{record}(a).\]
\lean{KIP126.Computation.Standard.atom_record}
\end{theorem}

\begin{theorem}[The h1 multiple of the correction term]
\label{frontier:E2Computation:6b059d10e824}
\level{fine}
For $C=e_0\Delta h_6g$, the sphere $E_2$-product satisfies
\[h_1C=0\in E_2^{14,139}(S^0).\]
\lean{KIP126.Computation.Standard.h1_correction_zero}
\end{theorem}

\begin{theorem}[Nonvanishing of the Kervaire square on E2]
\label{frontier:E2Computation:8116dddd2cea}
\level{fine}
The designated square satisfies
\[h_6^2\ne0\in E_2^{2,128}(S^0).\]
\lean{KIP126.Computation.Standard.square_nonzero}
\end{theorem}

\begin{theorem}[The class T is not divisible by h0]
\label{frontier:E2Computation:9a558822725b}
\level{fine}
There is no $a\in E_2^{13,138}$ with $h_0a=T$.
\lean{KIP126.Computation.Standard.t_not_h0_multiple}
\end{theorem}

\begin{theorem}[The exhaustive second-line basis in stem 126]
\label{frontier:E2Computation:37381227a3a8}
\level{fine}
Every $x\in E_2^{2,128}$ in the sphere Adams spectral sequence satisfies $x=0$ or $x=h_6^2$.
\lean{KIP126.Computation.Standard.square_exhaustive}
\end{theorem}

\begin{theorem}[Nondivisibility by $h_2$ on E2]
\label{frontier:E2Computation:3a36ab1c8384}
\level{fine}
The following statements hold in the sphere Adams spectral sequence.
\begin{enumerate}
\item There is no $a\in E_2^{13,135}$ with $h_2a=T$.
\item There is no $a\in E_2^{10,132}$ with $h_2a=Y$.
\end{enumerate}
\lean{KIP126.Computation.Standard.t_not_h2_multiple,KIP126.Computation.Standard.y_not_h2_multiple}
\end{theorem}

\begin{theorem}[The top-cell comparison and a product relation]
\label{frontier:E2Computation:fd3a2ddf7954}
\level{fine}
The following statements hold.
\begin{enumerate}
\item Let $\delta:C\nu\to S^3[1]$ be the cofiber connecting map, and let $\sigma_{j,1}:S^j[1]\cong S^{j+1}$ be the fixed shift-coherence isomorphisms. There are classes
\[a_3\in E_2^{8,133}(S^3),\quad a_2\in E_2^{8,132}(S^2),\quad a_1\in E_2^{8,131}(S^1)\]
such that $\delta_*a$ desuspends to $a_3$, $(\sigma_{2,1}^{-1})_*a_3$ desuspends to $a_2$, $(\sigma_{1,1}^{-1})_*a_2$ desuspends to $a_1$, and $a_1$ desuspends to $X$. These are the four fixed sphere desuspension relations.
\item $Xh_2=0\in E_2^{9,134}(S^0)$.
\end{enumerate}
\lean{KIP126.Computation.Standard.top_comparison,KIP126.Computation.Standard.x_h2_zero}
\end{theorem}

\begin{theorem}[Source coordinates recover monomials and named atoms]
\label{frontier:E2Computation:d71314d20c81}
\level{fine}
For $s,t\in\mathbb N$, $t\leq261$ and $i<n_{s,t}$,
\[q_{s,t}^{-1}(e_i)=m_{s,t,i}(u),\]
where the right side evaluates the archive monomial using the fixed source unit, products and generator values. For every named atom $a$, transporting $u_{\iota(a)}$ along $\deg(\iota(a))=\deg(a)$ gives exactly $a^{\mathrm{src}}$.
\lean{KIP126.Computation.Standard.sourceBasis_values,KIP126.Computation.Standard.sourceAtom_agrees}
\end{theorem}

\begin{theorem}[Calibration of the six standard h classes]
\label{frontier:E2Computation:d76af9345531}
\level{fine}
For each $i\in\{0,1,2,4,5,6\}$, the source atom is sent to the designated standard sphere class:
\[c_{1,2^i}(h_i^{\mathrm{src}})=h_i\in E_2^{1,2^i}(S^0).\]
\lean{KIP126.Computation.Standard.comparison_h6,KIP126.Computation.Standard.comparison_h2,KIP126.Computation.Standard.comparison_h4,KIP126.Computation.Standard.comparison_h0,KIP126.Computation.Standard.comparison_h5,KIP126.Computation.Standard.comparison_h1}
\end{theorem}

\begin{theorem}[Complete bases of the nu-cofiber]
\label{frontier:E2Computation:6e2df2108f8d}
\level{fine}
For every degree record $(C\nu,s,t;m_0,\ldots,m_{n-1})$ in the fixed finite table, there is a linear equivalence
\[e:E_2^{s,t}(C\nu)\cong_{\mathbb Z}\bigoplus_{i<n}\mathbb F_2,\qquad e^{-1}(e_i)=b_{C\nu;s,t,i}.\]
\lean{KIP126.Computation.Table.cofiber_basis_correct}
\end{theorem}

\begin{theorem}[Commutativity of bounded sphere products]
\label{frontier:E2Computation:1652f2bec247}
\level{fine}
For $s,t,s',t'\in\mathbb N$ with $t+t'\leq261$, $x\in E_2^{s,t}(S^0)$ and $y\in E_2^{s',t'}(S^0)$,
\[xy=yx\in E_2^{s+s',t+t'}(S^0),\]
using the canonical identifications of the two sums of bidegrees.
\lean{KIP126.Computation.Table.product_comm}
\end{theorem}

\begin{theorem}[Evaluation of every recorded basis expression]
\label{frontier:E2Computation:724420370993}
\level{fine}
For every record $(e,i)$ in the fixed basis-value list, let $(s,t)$ be the bidegree of the expression $e$. Then
\[b_{S^0;s,t,i}=\operatorname{value}(e).\]
Here expression values use the fixed sphere unit, named atoms, addressed archive generators and sphere products.
\lean{KIP126.Computation.Table.basis_value}
\end{theorem}
